\chapter{Newborn Hearing Screening}

\section*{Overview}
Newborn hearing screening is a quick, painless screening procedure performed soon after birth to identify babies who may have hearing loss.

\section*{How It Is Performed}
While the baby is asleep or calm, the clinician places a small probe in the ear for an otoacoustic-emissions test or soft earphones and surface sensors for an automated auditory brainstem-response test. The equipment records the ear's or auditory pathway's response to sound.

\section*{Why It Is Important}
Early identification allows confirmatory testing and timely support for speech, language, and learning. A screening result is not a diagnosis of permanent hearing loss.

\section*{Preparation and Results}
No special preparation is needed. The baby should be comfortable and quiet. A result of ``pass'' means no further immediate screening is usually needed, while a ``refer'' result requires repeat screening or diagnostic evaluation. Fluid or debris in the ear and movement can affect the result.

\section*{Risks}
The procedure is noninvasive and does not cause pain. The main limitation is that screening can occasionally miss hearing loss or produce a false-positive result.

\section*{References}
\begin{enumerate}
  \item National Institute on Deafness and Other Communication Disorders. Your baby's hearing and communicative development checklist. 2025.
  \item Current Medical Diagnosis and Treatment. McGraw-Hill; 2025.
\end{enumerate}
\clearpage
